\section{The center observable}\label{sec:observable}

We determine the change from \(M_{L^N}\) to \(M_{L^{N+1}}\).
The two computations use the same entry
scale and the same stopping scale.  We choose their local maps so that
the field-independent factors extracted near the insertion agree
exactly.  After controlling the surviving labelled term, we compare
the remaining terminal integrals and read the gradient coefficient
of Section~\ref{sec:trajectory}.

Write \(\eta_n\) for the minimizing unit field in the free dual square
corresponding to \(M_n\), as in Lemma~\ref{ent:center-background}, and
put
\[
                         g_n=\|\eta_n\|^2.
\]
All energies in this section count each unoriented edge once.

\begin{proposition}[Consecutive square sizes]\label{obs:consecutive}
Fix \(0<\epsilon<1/10\).  For all sufficiently large \(N\), put
\(n_0=L^N\), \(n_1=L^{N+1}\), and choose a power \(m'=L^{j'}\) so that
\begin{equation}\label{obs:stop}
 D_N=\frac{n_0}{m'}\asymp N^\epsilon,\qquad
 j'=N-O(\log N).
\end{equation}
Use the entry history for \(N\) in both squares, and let \(t_{j'}\)
be its unbalanced plane gradient coefficient at the stop.  Then
\begin{equation}\label{obs:consecutive-formula}
 \log\frac{M_{n_1}}{M_{n_0}}
 =-\frac{1-t_{j'}}{2a}(g_{n_1}-g_{n_0})
   +O\left(\frac{D_N^4}{N^2}+\frac{1}{ND_N}\right).
\end{equation}
Moreover, for a constant \(c_G\),
\begin{equation}\label{obs:green-constant}
                g_{L^N}=2\pi(\log L)N+c_G+o(1).
\end{equation}
All constants are independent of \(N\).
\end{proposition}

The term \(1-t_{j'}\) separates the Gaussian cost of the unit
circulation from the correction read from the height interaction.
The displayed error is summable for the chosen growth of \(D_N\).
We first prove the reference-energy assertion
\eqref{obs:green-constant}, and then prove the consecutive-size
comparison.

\paragraph{The reference energy.}
The unit field is \(2\pi\) times the rotated gradient of the killed
primal Green function \(G_n(\,\cdot\,,0)\).  The stream identity gives
\begin{equation}\label{obs:stream-energy}
 g_n=4\pi^2\langle G_n(\,\cdot\,,0),
                 A_n^{\rm kill}G_n(\,\cdot\,,0)\rangle
     =4\pi^2G_n(0,0).
\end{equation}
For the sequence \(n=L^N\), choose one sufficiently large fixed dyadic
\(D_G\), distinct from the growing \(D_N\), and put \(v=n/D_G\).
For late \(N\), \(v\) is a dyadic integer.  Write
\(\mathsf P_t^{\rm kill}=P_t(1-A_n^{\rm kill}/64)\) for the
polynomials of \eqref{ent:reference-split}.  Since
\(\mathsf P_1^{\rm kill}=I\), their exact telescope is
\begin{equation}\label{obs:green-telescope}
 (A_n^{\rm kill})^{-1}
 =\sum_{\substack{t=1,2,4,\ldots\\t<v}}
     (A_n^{\rm kill})^{-1}
       (\mathsf P_t^{\rm kill}-\mathsf P_{2t}^{\rm kill})
       +(A_n^{\rm kill})^{-1}\mathsf P_v^{\rm kill}.
\end{equation}
There is no independent microscopic piece in this identity.
The quotients use their polynomial continuations when compared with
the plane operators.

Choose \(D_G\) larger than the fixed range constant.  Finite range
then makes every center diagonal in the sum in
\eqref{obs:green-telescope} equal to its plane diagonal
\[
 H_t(0,0),\qquad H_t=A^{-1}(\mathsf P_t-\mathsf P_{2t})
\]
for the plane nearest-neighbor Laplacian.  The fixed-ratio Fourier
estimate used in the proof of Lemma~\ref{trj:coefficient-rate}
applies at ratio \(2\).  To check the omitted microscopic factor in
\cite[Lemma~8.1]{B}, its nearest-neighbor kernel contains
\(1-\gamma Q\), with \(Q=A/4\).  The difference from \(H_t\) is
\(\frac{\gamma}{4}(\mathsf P_t-\mathsf P_{2t})\); its plane diagonal
is \(O(t^{-2})\), by
\(\widehat{\mathsf P_t}(\xi)\le C(1+t|\xi|)^{-64}\).
This correction and the Fourier power error are summable over dyadic
\(t\).  The continuum diagonal is \((\log 2)/(2\pi)\): if \(P(r)\)
is the limiting radial cutoff, then
\[
 \int_0^\infty\bigl(P(r)-P(2r)\bigr)\,\frac{\dd r}{r}=\log 2
\]
because \(P(0)=1\) and \(P\) decays at infinity, and the angular
factor is \(1/(2\pi)\).  Thus, including the finitely many early sides,
\[
 H_t(0,0)=\frac{\log 2}{2\pi}+\varepsilon_t,\qquad
             \sum_{t=1,2,4,\ldots}|\varepsilon_t|<\infty.
\]
The sum of the center diagonals in \eqref{obs:green-telescope} is
therefore \((2\pi)^{-1}\log v\) plus a convergent constant and \(o(1)\).

The last center diagonal in \eqref{obs:green-telescope} has a limit
for this fixed \(D_G\).  For the normalized positive sine modes
\(e_{k,n}\) and their Dirichlet eigenvalues \(\lambda_{k,n}\),
\[
 |e_{k,n}(0)|^2\le Cn^{-2},\qquad
 \lambda_{k,n}\ge c|k|^2n^{-2},\qquad
 \mathsf P_v^{\rm kill}(k)
       \le C\min\{1,(|k|/D_G)^{-64}\}.
\]
Its summands are consequently bounded by
\(C|k|^{-2}\min\{1,(|k|/D_G)^{-64}\}\), a summable function of the
positive mode indices when \(D_G\) is fixed.  Each fixed mode
converges with its Dirichlet normalization, so dominated convergence
gives the limit.  Since \(\log v=\log n-\log D_G\),
\[
 G_{L^N}(0,0)=\frac{\log L}{2\pi}N+c+o(1).
\]
Together with \eqref{obs:stream-energy}, this proves
\eqref{obs:green-constant}.  No rate is needed for its convergent
remainder.

\subsection{The paired local computation}

Take the entry side \(m=L^{j_0}=qRs\) supplied by
Theorem~\ref{ent:initialization} for the epoch containing \(N\).
Thus \(s\asymp N\) and \(j_0=O(\log N)\).  Use this same choice in
the two squares, with grids centered on the same unit.  At absolute
index \(j\), put \(m_j=L^j\).  All steps have \(j_0\le j<j'\), so
\[
                    m_{j+1}\le m'=\frac{n_0}{D_N}.
\]
The centered and unit sectors within a square have the same Gaussian
covariance; their difference is in the activities and in the factor
\(e^{-g_n/(2a)}\) in \eqref{eq:noisy-reference}.

Fix the finite set of regulator and weight reserves allowed by
Theorem~\ref{ent:initialization}.  Use a stronger norm for the
untwisted activities than for their translated unit copies, and
leave the weakest block weight, denoted \(A_w\), large enough for the
rooted sums below.  These choices are fixed before \(N\).  In a
recurrence, \(\|\cdot\|_j\) denotes the chosen norm for that class of
activities at side \(m_j\), with the same choice at successive
indices.  Define the size of the common plane bulk state by
\[
          U_j=|t_j|+|z_j|+\|\mathcal R_j\|_j.
\]
Lemma~\ref{trj:vanishing} keeps \(U_j\) uniformly small throughout
the history, and Proposition~\ref{trj:trajectory} gives \(U_j=O(N^{-1})\)
on every fixed-fraction interval ending at \(j'\).  Such an interval
is always taken with a slightly longer interval behind it in the
same history.

We now specify the exact algebra whose scalars will cancel.
For the free dual square \(\mathcal G_n\), the entry theorem gives
\[
 \mathcal H_{\mathcal G_n,m}^0(\phi)
       =c_{\mathcal G_n,m}\cZ_m(K_m^{0,n};\phi),\qquad
 \mathcal H_{\mathcal G_n,m}^\tau(\phi)
       =c_{\mathcal G_n,m}[x]\cZ_m(K_m^{\tau,n}+xJ_m^n;\phi).
\]
Here \(K_m^{\tau,n}\) sums unit-sector components without the
compulsory source marker, while \(J_m^n\) sums the component containing
that marker.  Thus \([x]\) selects exactly one labelled component.
The positive entry scalar is common to the two sectors of this
square; it need not be common to the two square sizes.

Use the exact unbalanced partition map and work in
\(\R[x]/(x^2)\).  Let \(\mathcal B_{m_{j+1}}^n\) be the output
block grid in square \(n\), and let \(\zeta_j^n\) have the shell
covariance \(\Gamma_{m_j}\) there.  The unit step has the exact form
\begin{equation}\label{obs:dual-number-map}
 \begin{split}
 &\E_{\zeta_j^n}
    \cZ_{m_j}(K_j^{\tau,n}+xJ_j^n;\phi+\zeta_j^n)\\
 &\quad =
 \prod_{B\in\mathcal B_{m_{j+1}}^n}
       (1+a_{jB}^{\tau,n}+x\dot a_{jB}^{\tau,n})\,
 \cZ_{m_{j+1}}(K_{j+1}^{\tau,n}+xJ_{j+1}^n;\phi)
 \pmod{x^2}.
 \end{split}
\end{equation}
Use \(a_{jB}^{0,n}\) for the corresponding centered scalar.
Each \(a\) is the complete sum of designated field-independent pieces
assigned to \(B\); \(\dot a\) is its coefficient in the label
direction.  Only the unlabelled \(a\) must be small, since
\[
 (1+a+x\dot a)^{-1}
       =(1+a)^{-1}-x\dot a(1+a)^{-2}\pmod{x^2}.
\]
The nonlinear map is a convergent sum of bounded multilinear maps
of degree at least two, so
\[
                    \|D\mathcal N(K)J\|
                        \le C\|K\|\,\|J\|
\]
at a small unlabelled input, for every finite-norm direction \(J\).
No small-ball condition on \(xJ\) at a fixed nonzero \(x\) is used.

There are two copying operations.  Within a square, an activity away
from the source is the untwisted activity at a local lift.  We keep
its translated nonconstant pieces inside the activity and designate
the untwisted constant; this removes scalar differences away from the
source.  Between the two squares, central unit activities are the same
plane-unit prescription at different smooth backgrounds.  Designating
the neutral constant before that background is inserted makes the
remaining scalar independent of the square size.  The central
activities themselves retain the background difference.

For the precise support conditions, include in the causal padding
\(D_{m_j}(X)\), its fixed stencils, and an additional collar of fixed
positive width in current-block units.  Call a support \emph{wall-free}
when its required sites and causal padding stay a fixed number \(H_0\)
of current blocks from the square boundary.  Call it \emph{shift-good}
when its enclosing rectangle with that padding misses the unit; a
support failing this condition is a \emph{defect}.  The full norm
domain of a shift-good support is then at distance at least \(c m_j\)
from the unit.  Choose \(H_0\) once to include the shell range,
assignment neighborhoods, and the geometric sum of all earlier causal
padding.  Occupied support,
norm padding, and causal dependence have distinct roles here: every
obstruction tag remains occupied in the support; the norm padding
contains the evaluated coordinates; the causal neighborhood contains
the predecessor data used by the local operations.  The entry copy
conditions are those of Theorem~\ref{ent:initialization}.
These classifications are used in the proofs of
\cite[Propositions~5.5 and~5.7]{B}, whose stated results concern a
decaying Villain flow.  We derive the needed estimates for the
critical height history.

Fix a constant \(H_A\) large enough for the assignments described
below.  In the common centered grid define the geometric set
\begin{equation}\label{obs:affected-set}
 \mathcal A_j=
 \{\,m_{j+1}\text{-blocks meeting the }H_A m_{j+1}
                         \text{-neighborhood of the unit}\,\}.
\end{equation}
This is a set of possible scalar destinations, not the set where a
numerical difference happens to be nonzero.  Its cardinality is
bounded by a fixed constant.  For late \(N\), it lies in both squares
because \(m_{j+1}\le m'\ll n_0\).

Let \(w_j\) be the maximum over the two squares of the full untwisted
norm restricted to supports that are not wall-free, and let \(d_j\)
be the corresponding maximum of the full unlabelled unit norm on
defects.  Write \(\|J_j\|_j=\max_{n\in\{n_0,n_1\}}\|J_j^n\|_j\).
Each norm here is the fixed choice for its class.

\begin{lemma}[Paired local maps]\label{obs:local-estimates}
The local maps can be chosen in both squares with the following
properties.  Here \(\theta<1\) and \(C\) depend only on the fixed
parameters.
\begin{enumerate}
\item Untwisted activities on wall-free supports are the exact plane
activities, and
\begin{equation}\label{obs:wall-recurrence}
 w_{j+1}\le \theta w_j+C U_j+C(U_j+w_j)^2.
\end{equation}
\item On a shift-good support, \(K_j^{\tau,n}\) is the corresponding
untwisted activity evaluated at a local lift \(\lambda_n\), with
\(\nabla\lambda_n=\eta_n/\alpha\).  Near a wall this uses the
finite-square untwisted activity.  On defects,
\begin{equation}\label{obs:defect-recurrence}
 d_{j+1}\le\theta d_j+C(U_j+w_j)+C(U_j+w_j+d_j)^2.
\end{equation}
On a central support whose complete causal neighborhood is wall-free
in both squares, the unit activity is the common plane centered-unit
prescription evaluated at
\begin{equation}\label{obs:smooth-background}
                     \phi+\delta_n,\qquad
          \delta_n=\frac{v_n-v_\infty}{\alpha}.
\end{equation}
\item The surviving label obeys
\begin{equation}\label{obs:label-recurrence}
 \|J_{j+1}\|_{j+1}
 \le \{\theta+C(U_j+w_j+d_j)\}\|J_j\|_j.
\end{equation}
Every nonzero \(J_j^n(X)\) has an occupied block within a fixed number
of current blocks of the unit.  Its support may extend farther.
It has the same central plane-unit copy property as in
\eqref{obs:smooth-background}.
\item Choose \(H_A\) in \eqref{obs:affected-set} sufficiently large.
For every output block outside \(\mathcal A_j\), in either square,
\begin{equation}\label{obs:factor-outside}
 a_{jB}^{\tau,n}=a_{jB}^{0,n},\qquad
                  \dot a_{jB}^{\tau,n}=0.
\end{equation}
For every \(B\in\mathcal A_j\), the complete sums satisfy
\begin{equation}\label{obs:factor-copy}
 a_{jB}^{0,n_0}=a_{jB}^{0,n_1},\qquad
 a_{jB}^{\tau,n_0}=a_{jB}^{\tau,n_1},\qquad
 \dot a_{jB}^{\tau,n_0}=\dot a_{jB}^{\tau,n_1}.
\end{equation}
In the fixed small ball, choose \(\rho<1\) with
\(|a_{jB}^{\sigma,n}|\le\rho\) for \(\sigma\in\{0,\tau\}\).
Then every output block satisfies
\begin{equation}\label{obs:scalar-log-bound}
 \left|\log\frac{1+a_{jB}^{\tau,n}}{1+a_{jB}^{0,n}}\right|
       \le C(U_j+w_j+d_j).
\end{equation}
\end{enumerate}
\end{lemma}

\begin{proof}
We first establish the copying and scalar identities at a step for
which the unlabelled inputs lie in the small ball.  The estimates
below keep the iterated maps in that ball.

On a shift-good input use its full translated untwisted expectation.
Translate also the already assigned fine-site gradient and fundamental
functions.  Thus both \(e_B^0(\phi+\lambda_n)\) and
\(c_B^\alpha(\phi+\lambda_n)\), including their values at zero, remain
functions inside the interaction.  Designate only the same
field-independent constant as in the untwisted map.  These choices
are exact splittings in \cite[Proposition~4.8]{B}.  The weak map
assigns its fundamental functions to the fine blocks \(B\in X\)
before a support label is moved
\cite[Section~8.2, before Proposition~8.3]{B}.  A piece that later
belongs to a defect output is still that fine-site function, so no
lift is evaluated at the unit.  Lifts on opposite sides of the cut
differ by a common integer multiple of \(\omega\); periodicity
identifies their resulting functions.

On a small defect input, designate the neutral constant of the plane
centered-unit prescription before the background \(\delta_n\) is
included.  This also applies coefficientwise to a labelled input.
In square coordinates it is the neutral shell expectation evaluated
at \(-\delta_n\).  The initial central copy identity, including its
occupied tags and padded dependencies, is in
Theorem~\ref{ent:initialization}.  In a central causal neighborhood
the shell covariance equals the plane covariance.  Translation of
the entire field argument commutes with convolution and products,
and the prescribed neutral constant is the plane constant.  Induction
therefore gives \eqref{obs:smooth-background}, also coefficientwise
in \(x\).  The same induction gives the untwisted plane copy on
wall-free supports.  Near a wall a shift-good input is instead
compared to its finite-square untwisted input and its assigned pieces.
The iterated causal radius is \(O(m_j)\), since earlier radii form
a geometric sum.

We next locate every possible scalar difference.  Here a small
input means \(|X|\le n_{\rm sm}\) for its full occupied support
\(X\), with every obstruction tag still occupied; the norm padding
and causal neighborhood enter the support predicates just defined.
For a current small linear input \(X\), write
\(c_{jX,B}^{\sigma,n}\) for the complete designated scalar assigned
from that input to \(B\) in its coarse closure \(\bar X\).
The exact algebra of \cite[Proposition~4.8]{B} gives
\begin{equation}\label{obs:scalar-sums}
 a_{jB}^{\sigma,n}
       =\sum_{\substack{X\ {\rm small}\\B\in\bar X}}
                      c_{jX,B}^{\sigma,n},\qquad
 \dot a_{jB}^{\tau,n}
       =\sum_{\substack{X\ {\rm small}\\B\in\bar X}}
                      \dot c_{jX,B}^{\tau,n}.
\end{equation}
The second sum is the same linear designation in the label direction.
A large or nonlinear residual is not designated at that step, even
if its value happens to be constant; it retains its whole coarse
closure.  It can enter \eqref{obs:scalar-sums} at a later step only
after its then-current full occupied support is small.

The shift-good designation is exactly the centered one.  Within
one square, therefore,
\begin{equation}\label{obs:scalar-difference}
 \Delta a_{jB}^n:=a_{jB}^{\tau,n}-a_{jB}^{0,n}
   =\sum_{\substack{X\ {\rm small\ defect}\\B\in\bar X}}
             (c_{jX,B}^{\tau,n}-c_{jX,B}^{0,n}).
\end{equation}
Every nonzero summand has a small defect witness.  Its fixed connected
size and its padded rectangle meeting the unit put its entire support
and destination within \(O(m_{j+1})\) of the unit.  This argument
does not assign source ancestry to an arbitrary unlabelled defect:
a large winding support can enclose the unit without occupying a
nearby block.

A labelled input has a different witness.  Initially its occupied
cover contains the source by Theorem~\ref{ent:initialization}.  A
surviving label comes from its fine input or from a labelled
normalization decoration.  A piece can shrink only when it came
from a small linear input, and such a piece or decoration moves by
only \(O(1)\) output blocks.  The accumulated physical displacement
is \(O(m_j)\).  Thus every surviving labelled support has an occupied
block within a fixed number of current blocks of the unit.  A small
labelled input and its scalar destination are consequently within
\(O(m_{j+1})\) of the unit.  Choosing \(H_A\) for both kinds of
witness proves \eqref{obs:factor-outside}.

We must also rule out noncommon numerical data hidden inside a
complete scalar at \(B\in\mathcal A_j\).  For the inverse decoration
\(\mathfrak d_{jB}^{\sigma,n}=(1+a_{jB}^{\sigma,n})^{-1}-1\), the
exact difference and labelled derivative are
\[
 \mathfrak d_{jB}^{\tau,n}-\mathfrak d_{jB}^{0,n}
   =-\frac{\Delta a_{jB}^n}
           {(1+a_{jB}^{\tau,n})(1+a_{jB}^{0,n})},
 \qquad
 \dot{\mathfrak d}_{jB}^{\tau,n}
   =-\frac{\dot a_{jB}^{\tau,n}}{(1+a_{jB}^{\tau,n})^2}.
\]
Every other factor in either inverse series is a complete scalar at
this same block.  A decoration at a remote block occupies that block
when grouping adds it to the support.  Large residuals retain their
closures; a permitted small relocation keeps its old exclusion
information and moves a witness only \(O(m_{j+1})\).  In units of
the new side its distance bound has the form
\(h_{\rm new}\le C+h/L\), so repeated allowed relocations retain a
fixed bound.

The entry construction occupies every wall-reaching or winding
obstruction tag.  If a current unlabelled defect has a wall witness
and its padded rectangle meets the unit, its occupied extrema in a
wall-normal coordinate are separated by \(n-O(m_j)\).  Fixed-distance
connectivity then gives at least \(c n/m_j\) occupied blocks.
The same span bound applies when a support joins an affected inverse
witness to a wall-obstruction tag.  Every noncommon finite-square
coefficient correction carries such a tag.  Through the stop this
lower bound is at least \(cD_N\), so
such a support cannot become a small central input.  This is why a
long span cannot be erased before the smallness decision.

At a block \(B\in\mathcal A_j\), every contributor to the complete
sums in \eqref{obs:scalar-sums} is itself small and lies within
\(O(m_{j+1})\) of \(B\).  Its locally evaluated coordinates and rule
data are therefore identified in the common central region by the causal padding.
A noncommon finite-square coefficient correction would carry a
wall-obstruction tag, contradicting this small central support by the
preceding span bound.  All remaining untagged coefficients use the
same universal plane prescription.  The centered contributions
therefore copy term by term between sizes.  For a central unit input
written as
\(H_n(\phi)=H_\infty(\phi+\delta_n)\), its designated neutral
constant is \(H_n(-\delta_n)=H_\infty(0)\); a shift-good unit input
designates the centered constant.  The same reasoning holds in the
label direction.  This proves \eqref{obs:factor-copy} for the
complete centered sum, complete unit sum, and complete labelled
derivative separately.  It is stronger than equality merely of
\(\Delta a_{jB}^n\), which would not identify their ratios.

We turn to the analytic estimates.  At a free wall, a small
untwisted input is invariant under field negation.  Removing its
neutral constant gains \(CL^{-2}\) by \cite[Lemma~4.5]{B}, and there
are \(O(L)\) fine placements at a specified coarse wall block by
\cite[Corollary~4.7]{B}.  The fundamental charged part has the same
\(CL^{-2}\) gain before placements at the present marginal frequency.
Indeed the plane shell diagonal is
\[
                         \frac{\log L}{2\pi}+O(1).
\]
For a reflected image at distance \(d\), binary pieces of side less
than a fixed multiple of \(d\) vanish.  A later piece of side \(t\)
differs from its nonnegative zero-distance value by at most \(Cd/t\),
using the first-difference kernel bound.  Summing over dyadic
\(t\ge cd\) costs \(O(1)\), and only a fixed number of images occur
through the stop.  The free diagonal is therefore at least
\((\log L)/(2\pi)-C\).  For each small connected charged input \(X\),
choose the local anchor
\(x=v_0\) in one of its occupied blocks, and apply the general charge
estimate \cite[Equation~(4.26)]{B} with Cameron shift amplitude
\(\alpha\).  For charge \(q\) its diagonal exponent is
\(-\alpha^2(|q|-1/2)\Gamma(x,x)\).  Since \(\alpha^2=8\pi\), this
gives \(CL^{-2}\) for \(|q|=1\), and a summable bound with spare
powers of \(L\) for the other charges.  The fixed analytic radius
contains the nonconstant covariance column; choose \(L\) after
this radius.

Large linear supports contract by the block weight.  A bulk input
feeding a wall output, including every piece assigned from it, is
bounded in full by \(CU_j\); this uses no cancellation across a
change of boundary kernel.  The analytic nonlinear estimate
\cite[Equation~(4.30)]{B} now gives \eqref{obs:wall-recurrence}.

For a shift-good input, \eqref{ent:unit-decay} bounds the first two
differences of its lift on each padded block after multiplication
by the corresponding powers of \(m_j\).  Its translated ordinary
regulator is bounded by a slightly stronger regulator times
\(e^{C|X|}\).  For the observation regulator, its positive quadratic
form, the strict margin in \eqref{eq:observation-regulator}, and
\cite[Equations~(9.14) and~(9.16)]{B} give the same conclusion:
the internal energy of the lift is \(O(|X|)\), and a fixed increase
of \(h_*\) absorbs the cross term.  The weight reserve pays
\(e^{C|X|}\).  Common lift values cost nothing in the regulators
or in the supremum over real backgrounds in the activity norm.
Contributions from shift-good inputs to a defect output are
therefore bounded by \(C(U_j+w_j)\).

The plane-before-\(\delta_n\) projection on a small defect is close
to ordinary neutral constant localization.  After removing a common
constant,
\[
                    \|\delta_n\|_{m_j,D}\le C m_j/n
\]
by \eqref{ent:unit-decay}.  The mean value theorem in the analytic
point norm, with the shell regulator reserve, bounds the additional
linear remainder by \(C_{L,A_w}(m_j/n)\) times the input norm.
This is at most \(C_{L,A_w}/D_N\).  Ordinary neutral constant
localization without negation symmetry gains \(CL^{-1}\), and only
\(O(1)\) fine defect placements can feed a specified coarse block.
The charged estimate above, with the local anchor in each small input,
and the large-support weight bound handle the other linear terms.  The
chosen projection has bounded
norm, so the nonlinear bound is still quadratic.  Taking \(N\)
large after the fixed map choices absorbs \(C_{L,A_w}/D_N\) in
the contraction and proves \eqref{obs:defect-recurrence}.

The linear label map uses this same neutral projection, the
order-one neutral remainder and charged gain on a small support with
its own local anchor, and the weight gain on a large support.  The
occupied-block
location proved above gives \(O(1)\) small labelled placements
at each step.  The differentiated nonlinear estimate following
\eqref{obs:dual-number-map} accounts for all ordinary attachments
and proves \eqref{obs:label-recurrence} for an arbitrary finite
label norm.

Finally, the scalar projections in \eqref{obs:scalar-sums} are
bounded in each chosen class norm, and the rooted sum over the
fixed-threshold small inputs assigned to one block is finite.
In \eqref{obs:scalar-difference}, the unit defect projections cost
\(Cd_j\).  The plane-before-\(\delta_n\) evaluation remains uniformly
bounded by the preceding projection estimate.  Bound the centered
counterparts separately in their own norm by \(C(U_j+w_j)\).
No difference of norms and no norm of \(J_j\) is used.  Hence
\[
                |\Delta a_{jB}^n|\le C(U_j+w_j+d_j).
\]
On \(|a_{jB}^{\sigma,n}|\le\rho<1\), the real mean value theorem
for \(\log(1+x)\) gives
\[
 \left|\log\frac{1+a_{jB}^{\tau,n}}{1+a_{jB}^{0,n}}\right|
 \le \frac{|\Delta a_{jB}^n|}{1-\rho}.
\]
This proves \eqref{obs:scalar-log-bound} and completes the lemma.
\end{proof}

Choose the entry and bulk tube small enough that the coefficient
in \eqref{obs:label-recurrence} is at most a fixed \(\theta'<1\).
The recurrences then keep the actual unlabelled maps in the required
small ball.  The polynomial entry bound gives
\begin{equation}\label{obs:label-stop}
 \|J_{j'}\|_{j'}\le C_Rs^{C_R}(\theta')^{j'-j_0}
                       \le e^{-cN}.
\end{equation}
The constant \(c>0\) may decrease between occurrences.  Iterating
\eqref{obs:wall-recurrence} and \eqref{obs:defect-recurrence} from
an earlier fixed-fraction interval, with the quadratic small
factors absorbed in the contraction, gives
\begin{equation}\label{obs:unlabelled-stop}
        U_{j'}+w_{j'}+d_{j'}\le C/N.
\end{equation}
The same bound holds for the plane centered-unit defect prescriptions.
More generally, Lemma~\ref{trj:vanishing} and these geometric
convolutions show that \(w_j+d_j=o(1)\) uniformly on every
fixed-fraction interior interval.  Earlier portions of the history
contribute exponentially little.

\subsection{The stopped identity and terminal integration}

At the stop, write
\[
 C_{*,n}=A_n^+\mathsf P_{m'}^n,\qquad
 \mathsf P_t^n=P_t(1-A_n/64),\qquad
 \Cov(\zeta_{*,n})=C_{*,n}
\]
where \(A_n\) is the free-square Laplacian and \(A_n^+\) is its
inverse on the nonconstant modes.  This is the terminal covariance in \eqref{ent:reference-split}.  The remaining common
phase is uniform modulo \(\omega\).  For \(\sigma\in\{0,\tau\}\),
define the three terminal quantities
\begin{equation}\label{obs:terminal-quantities}
\begin{split}
 T_n^\sigma
   &=\frac1\omega\int_0^\omega
       \E_{\zeta_{*,n}}
       \cZ_{m'}(K_{j'}^{\sigma,n};\zeta_{*,n}+t)\dd t,\\
 \dot T_n^\tau
   &=[x]\frac1\omega\int_0^\omega
       \E_{\zeta_{*,n}}
       \cZ_{m'}(K_{j'}^{\tau,n}+xJ_{j'}^n;
                                  \zeta_{*,n}+t)\dd t.
\end{split}
\end{equation}
Thus \(T_n^\tau\) contains no compulsory component, whereas
\(\dot T_n^\tau\) contains one surviving labelled component.

Let \(\widehat M_n\) denote the insertion ratio with the observation
cutoff included in both physical sectors.  Iterating
\eqref{obs:dual-number-map}, integrating the terminal field, and
taking the coefficient of \(x\) give the exact stopped identity
\begin{equation}\label{obs:extraction-identity}
 \widehat M_n=e^{-g_n/(2a)}D_{\rm ex}
              \frac{c_{\rm ex}T_n^\tau+\dot T_n^\tau}{T_n^0}.
\end{equation}
The common entry scalar cancels within this square.  Because the
actual local maps remain in their small ball, the outside cancellation
and complete-scalar copying in Lemma~\ref{obs:local-estimates} allow
us to write the remaining expressions exactly as
\begin{equation}\label{obs:extracted-scalars}
 D_{\rm ex}=
 \prod_{\substack{j_0\le j<j'\\B\in\mathcal A_j}}
          \frac{1+a_{jB}^{\tau,n}}{1+a_{jB}^{0,n}},
 \qquad
 c_{\rm ex}=
 \sum_{\substack{j_0\le j<j'\\B\in\mathcal A_j}}
          \frac{\dot a_{jB}^{\tau,n}}{1+a_{jB}^{\tau,n}}.
\end{equation}
Both are independent of \(n\in\{n_0,n_1\}\), although they may
depend on the pair.  Each factor uses the complete scalar assigned
to its block.  All its unlabelled denominators are positive, and
\(D_{\rm ex}>0\).  We next estimate the terminal quantities in
\eqref{obs:extraction-identity} and then justify discarding its last
term relatively.

The first estimate is uniform as the number of terminal blocks grows.
It uses the same observation regulator as the local norms.

\begin{lemma}[Terminal regulator bound]\label{ent:terminal-regulator}
With strict reserve in the parameters of
\eqref{eq:observation-regulator} (\(1<p_1<p_2\) and
\(0<h_*<1/p_2\)), let \(u\) be a later block side compatible with
the fixed entry observation grid.  Let the square or torus and its
\(u\)-block grid be aligned and compatible, with side \(M\ge2u\).
Then, with the terminal covariance \(C_{\ge u}\) from
\eqref{ent:reference-split}, for every compatible padded union \(X\)
of \(u\)-blocks and sufficiently small \(\kappa\),
\[
 \E\bigl[W_u^\kappa(X,\zeta)e^{V_u(X,\zeta)}\bigr]
 \le C^{|X|},\qquad \Cov(\zeta)=C_{\ge u}.
\]
The terminal Gaussian is on nonconstant modes; the uniform common
mean is integrated separately modulo \(\omega\).  This bound is uniform
as \(M/u\) increases.
\end{lemma}

\begin{proof}
We first prove a positive-order trace estimate.  It uses only integral
\(u\) and \(M\ge2u\), independently of the observation-cell grid,
so it also applies at the binary tail side used below.  Stack the
differences
\[
 \bigl(u^{i-1}\nabla^i\zeta(x):1\le i\le4,\ x\in D_u(X)^+\bigr),
 \qquad \Cov(\zeta)=C_{\ge u},
\]
where the superscript \(+\) adds the fixed difference stencils.
For the nearest-neighbor symbol, the covariance multipliers of its
diagonal blocks are bounded on the fundamental frequency square by
\[
 C_i(u|\xi|)^{2i-2}(1+u|\xi|)^{-32}.
\]
Their suprema are bounded.  Their normalized frequency sums are
\(O(u^{-2})\) whenever \(M\ge2u\): count the \(O((M/u)^2)\) modes in
the disk of radius \(u^{-1}\), and sum the convergent outer frequency
annuli.  The excluded constant mode contributes nothing to a
positive-order difference.  The doubled-torus calculation gives the
same estimate after even reflection.  Positivity bounds the mixed
blocks by the diagonal ones.  The stacked covariance consequently
has bounded operator norm and trace at most \(C|X|\).

The local discrete Sobolev inequality in the proof of
\cite[Lemma~4.4]{B} bounds the exponent of \(W_u^\kappa\) by a
constant times the squared norm of this stack.  Its Gaussian
determinant is at most \(C^{|X|}\) for sufficiently small \(\kappa\).
For \(V_u\), let \(S_{\mathrm o}\) and \(S_{\mathrm t}\) be the maps
from the old Gaussian and noise coordinates and from the terminal
coordinates into the Hilbert space of the open observation energy.
Restricting one trial field to all
components of an open union can only lower its minimized observation
energy; each component may minimize its constant separately.  The
same joint trial restriction gives
\(S_{\mathrm o}S_{\mathrm o}^*+S_{\mathrm t}S_{\mathrm t}^*\le I\).
Put \(\beta=p_1h_*<1\) and \(A_{\mathrm o}=S_{\mathrm o}S_{\mathrm o}^*\).
Optimization of the old coordinates leaves, in the terminal Gaussian
integral of \(e^{p_1V_u}\), the positive matrix
\[
 \beta S_{\mathrm t}^*(I-\beta A_{\mathrm o})^{-1}S_{\mathrm t}.
\]
Its nonzero spectrum is that of
\(\beta(I-\beta A_{\mathrm o})^{-1/2}
S_{\mathrm t}S_{\mathrm t}^*(I-\beta A_{\mathrm o})^{-1/2}\).
The latter is bounded above by
\[
 \beta(I-\beta A_{\mathrm o})^{-1/2}
 (I-A_{\mathrm o})(I-\beta A_{\mathrm o})^{-1/2}
 \le \beta I,
\]
and its trace is at most
\(\beta(1-\beta)^{-1}\tr(S_{\mathrm t}S_{\mathrm t}^*)\).
This trace is bounded by
\[
 C\,\E E_{D_u(X)}^0(B_s\alpha\zeta)
 \le C\,\E\ip{\zeta}{A_{D_u(X)}\zeta}
 \le C|X|,
\]
using the same positive-order trace estimate.  The fixed spectral gap
therefore gives \(\E e^{p_1V_u}\le C^{|X|}\).  H\"older with the
reserved parameters combines the two determinants.  These arguments
use no bound on point-value variance; the uniform common mean is
integrated separately.
\end{proof}

\begin{lemma}[Terminal families and translated fields]\label{obs:terminal-family}
Write
\[
 W_u^\kappa(X,\phi)=\exp\{\kappa U_u^{\rm reg}(X,\phi)\}.
\]
Here \(U_u^{\rm reg}\) is the sum of gradient energy on \(D_u(X)\),
the \(u\)-weighted boundary trace, and the \(u^4\)-weighted block
maxima of squared second differences.  Its domain contains all
boundary endpoints and fixed difference stencils, with even extension
at a free wall.  Let
\[
 \mathfrak R_u^{\rm b}(X,\phi)
       =W_u^{\kappa_{\rm b}}(X,\phi)e^{V_{u,h_{\rm b}}(X,\phi)}
\]
be the regulator used to measure terminal activities, where
\(V_{u,h}\) denotes \eqref{eq:observation-regulator} with \(h_*=h\).

There is a fixed \(1<p_{\rm seg}<p_1\) with the following property.
For a support \(X\) of \(k\) terminal blocks, put \(D=D_{m'}(X)\).
Suppose a real deterministic field \(h_X\), allowed to depend on
\(n\) and \(X\), is defined on the full domain just described and
satisfies
\begin{equation}\label{obs:terminal-translation}
 U_{m'}^{\rm reg}(X,h_X)+\ip{h_X}{A_Dh_X}\le C_{\rm tr}k
\end{equation}
with one uniform constant.  Then
\begin{equation}\label{obs:terminal-moment}
 \E\bigl[\mathfrak R_{m'}^{\rm b}
                 (X,\zeta_{*,n}+h_X)^{p_{\rm seg}}\bigr]\le C^k.
\end{equation}
The same bound holds at every point of a coupled segment between
two such endpoints when their local regulator forms are identical.
A random common anchor may be subtracted at each endpoint after
neutral projection.  The assertion for the segment is a bound on
the moment at each parameter value, uniformly in that value.

There are \(H_N=O_L(D_N^2)\) terminal blocks in either square.
Suppose every unlabelled terminal activity has norm at most \(\nu\)
in the weakest norm and put \(b_N=CH_N\nu\).  If \(b_N=o(1)\), its
terminal partition integral \(T\) satisfies
\begin{equation}\label{obs:terminal-expansion}
 T=1+S_1+O(b_N^2),\qquad |S_1|\le b_N,\qquad
                  \log T=S_1+O(b_N^2),
\end{equation}
where \(S_1\) is the sum of single-activity expectations.
For a label direction of norm \(v\), the derivative obeys
\begin{equation}\label{obs:terminal-derivative}
                         |\dot T|\le CH_Nv e^{b_N}.
\end{equation}
These estimates allow the constituent-dependent translations in
\eqref{obs:terminal-translation}; their fixed per-block cost is
absorbed by the weight.
\end{lemma}

\begin{proof}
Choose a stronger admissible regulator
\(\mathfrak R_u^{\rm s}=W_u^{\kappa_{\rm s}}e^{V_{u,h_{\rm s}}}\)
from the fixed reserves, with
\(\kappa_{\rm s}>\kappa_{\rm b}\) and
\(h_{\rm b}<h_{\rm s}<1/p_2\).  Fix a small \(\varepsilon_{\rm tr}>0\) with
\((1+\varepsilon_{\rm tr})\kappa_{\rm b}\le\kappa_{\rm s}\) and
\((1+\varepsilon_{\rm tr})h_{\rm b}\le h_{\rm s}\).  The quadratic terms and maxima of squares give
\[
 U_u^{\rm reg}(X,\phi+h)
 \le (1+\varepsilon_{\rm tr})U_u^{\rm reg}(X,\phi)
          +(1+\varepsilon_{\rm tr}^{-1})U_u^{\rm reg}(X,h).
\]
Applying the same quadratic inequality inside the variational
definition of \(V\), while leaving its Cameron penalty unchanged,
gives
\[
\begin{split}
 V_{u,h_{\rm b}}(X,\phi+h)
 &\le V_{u,(1+\varepsilon_{\rm tr})h_{\rm b}}(X,\phi)
                         +C_{\varepsilon_{\rm tr}} E_D^0(B_s\alpha h)\\
 &\le V_{u,h_{\rm s}}(X,\phi)+C_{\varepsilon_{\rm tr}}\ip{h}{A_Dh}.
\end{split}
\]
The last inequality tries \(\alpha h\) in the open observation
minimum and uses \(\alpha^2=a\).  Therefore
\eqref{obs:terminal-translation} gives the pointwise comparison
\begin{equation}\label{obs:translated-regulator}
 \mathfrak R_u^{\rm b}(X,\phi+h_X)
                  \le e^{Ck}\mathfrak R_u^{\rm s}(X,\phi).
\end{equation}
The internal energy alone would not control the ordinary regulator:
its second-difference term also appears in
\eqref{obs:terminal-translation}.

Choose \(1<p_{\rm seg}<p_1\).  The matrix estimate in the proof of
Lemma~\ref{ent:terminal-regulator} gives
\(\E e^{p_1V_{m',h_{\rm s}}}\le C^k\), since
\(p_1h_{\rm s}<1\).  Reserve the ordinary determinant at exponent
\(p_{\rm seg}\kappa_{\rm s}p_1/(p_1-p_{\rm seg})\).
H\"older's inequality and \eqref{obs:translated-regulator} then
prove \eqref{obs:terminal-moment}.  For a compatible family with
different \(h_X\), apply \eqref{obs:translated-regulator} to each
constituent first.  The product inequality reduces the stronger
unshifted regulators to the regulator of their padded union.
Lemma~\ref{ent:terminal-regulator} bounds that union with cost
\(C^k\).  The terminal fields on its constituents need not be
independent.

The actual shifts in this section meet
\eqref{obs:terminal-translation}.  On a shift-good padded block,
\eqref{ent:unit-decay} bounds
\(u|\nabla\lambda_n|+u^2|\nabla^2\lambda_n|\) by a fixed constant.
On both the small wall patches and the central \(\delta_n\) patches
used below, the relevant shift \(h_X\) satisfies
\[
 u|\nabla h_X|\le C u/n,\qquad
 u^2|\nabla^2h_X|\le C(u/n)^2.
\]
There are \(O(ku^2)\) sites, \(O(ku)\) boundary
edges, and \(O(k)\) block maxima, including holes and fixed stencils.
Consequently \(U_u^{\rm reg}+\ip{h_X}{A_Dh_X}\le Ck\); on the
latter patches the ordinary cost is at most
\(Ck[(u/n)^2+(u/n)^4]\).  The regulator is unchanged by a common
constant.  Different lift branches differ by an integer multiple of
\(\omega\), so periodicity identifies their activity evaluations.

For the segment assertion, use the identical local forms at the
two endpoints.  The logarithm of \(\mathfrak R_u^{\rm b}\) is convex:
\(U_u^{\rm reg}\) is a sum of positive quadratic forms and maxima
of squares, and \(V_{u,h_{\rm b}}\) is a positive quadratic form.
Thus for any coupling and \(0\le r\le1\),
\[
 \mathfrak R_u^{\rm b}(X,(1-r)\xi_0+r\xi_1)
 \le \mathfrak R_u^{\rm b}(X,\xi_0)^{1-r}
       \mathfrak R_u^{\rm b}(X,\xi_1)^r.
\]
H\"older applies to the endpoint \(p_{\rm seg}\)-moments at each
\(r\).  A common anchor removed after neutral projection changes
neither the activity nor these regulator energies.  This proves the
stated uniform moment on the segment, without a second moment for
\(e^V\) or an expectation of a supremum over the segment.

For a compatible family with \(h\) constituents and total size
\(k\), the norm and the union bound now give the absolute integral
\[
                         \nu^h(C/A_w)^k.
\]
The rooted count \cite[Lemma~4.3]{B} sums one constituent to at most
\(CH_N\nu=b_N\).  Dropping compatibility enlarges an absolute sum,
so unordered families with \(h\) constituents cost at most
\(b_N^h/h!\).  Sum \(h\ge2\) and expand the logarithm to obtain
\eqref{obs:terminal-expansion}; in particular \(T>0\) for late
\(N\).  Distinguishing one labelled constituent gives \(CH_Nv\)
for it and \(e^{b_N}\) for all the others, proving
\eqref{obs:terminal-derivative}.
\end{proof}

Apply this lemma with \(\nu=C/N\) from
\eqref{obs:unlabelled-stop} and \(v\le e^{-cN}\) from
\eqref{obs:label-stop}.  Since \(D_N^2/N=o(1)\),
\begin{equation}\label{obs:terminal-near-one}
 T_n^0,T_n^\tau=1+O(D_N^2/N),\qquad
                 |\dot T_n^\tau|\le e^{-cN}.
\end{equation}

We need a lower bound on \(c_{\rm ex}\) only to discard the last
term of \eqref{obs:extraction-identity} relatively.  It follows
from the physical observable and requires no lower bound on the
label at entry.  The cutoff loss in \eqref{ent:cutoff-loss} and
\(M_n=n^{-1/8+o(1)}\) from Proposition~\ref{geo:comparison} give,
for the two current sizes,
\begin{equation}\label{obs:cutoff-relative}
                   \widehat M_n=M_n(1+O(e^{-cN^2})).
\end{equation}
Indeed the unit loss is absolute relative to the uncut centered
partition.  Dividing it by \(M_n=e^{-O(N)}\) still leaves the
displayed bound, and the centered loss has the same bound.

Set \(Q_n=\widehat M_n e^{g_n/(2a)}\).  The same rough center
asymptotic, \eqref{obs:cutoff-relative}, and
\eqref{obs:green-constant} imply
\[
                         Q_n=e^{o(N)}>0.
\]
The scalar estimate in Lemma~\ref{obs:local-estimates} also gives
\begin{equation}\label{obs:scalar-subexponential}
 |\log D_{\rm ex}|
   \le C\sum_{j=j_0}^{j'-1}(U_j+w_j+d_j)=o(N).
\end{equation}
For the last equality, there are only \(O(1)\) blocks in
\(\mathcal A_j\) at each index.  The summands are bounded in the
small ball and tend uniformly to zero on every fixed-fraction
interior interval.  For fixed \(\delta>0\), indices below
\(\delta N\) cost \(C\delta N+O(\log N)\); the rest cost \(o(N)\).
Let \(N\to\infty\) and then \(\delta\downarrow0\).

Rearranging \eqref{obs:extraction-identity},
\[
 c_{\rm ex}T_n^\tau
       =Q_n T_n^0/D_{\rm ex}-\dot T_n^\tau.
\]
The first term on the right is positive and at least \(e^{-o(N)}\)
by \eqref{obs:terminal-near-one} and
\eqref{obs:scalar-subexponential}; the second is exponentially
small.  Hence \(c_{\rm ex}>0\) for late \(N\),
\(c_{\rm ex}\ge e^{-o(N)}\), and
\(\dot T_n^\tau/(c_{\rm ex}T_n^\tau)=O(e^{-cN})\).
Taking logarithms in the two sizes now cancels both
\(D_{\rm ex}\) and \(c_{\rm ex}\) exactly.  With
\eqref{obs:cutoff-relative}, this gives
\begin{equation}\label{obs:after-cancellation}
 \log\frac{M_{n_1}}{M_{n_0}}
 =-\frac{g_{n_1}-g_{n_0}}{2a}
    +\left[\log(T_n^\tau/T_n^0)\right]_{n_0}^{n_1}
    +O(e^{-cN}),
\end{equation}
where \([F_n]_{n_0}^{n_1}=F_{n_1}-F_{n_0}\).

\subsection{The terminal energy and the other supports}

By Lemma~\ref{obs:terminal-family}, each terminal logarithm equals
its single-activity sum with error \(O(D_N^4/N^2)\).  We will prove
\begin{equation}\label{obs:terminal-energy}
 \left[\log(T_n^\tau/T_n^0)\right]_{n_0}^{n_1}
  =\frac{t_{j'}}{2a}(g_{n_1}-g_{n_0})
       +O\left(\frac{D_N^4}{N^2}+\frac{1}{ND_N}\right).
\end{equation}
The same \(t_{j'}\) occurs in both computations.

In a single-activity expectation, integration of the common phase
replaces an activity \(F\) by
\[
             F_0(\phi)=\frac1\omega\int_0^\omega F(\phi+t)\dd t.
\]
This neutral projection is invariant under every real common
constant and has the same norm bound
\cite[Section~4.3, before Lemma~4.5]{B}.  In particular the
fundamental singleton has zero expectation, including after a
nonconstant lift.

For a single-activity support \(X\), let \(k=|X|\) count its full
occupied blocks, including all obstruction tags, and put
\(D_X=D_{m'}(X)\).  Choose a fixed small \(c>0\) and use the following disjoint partition:
\begin{equation}\label{obs:terminal-partition}
\begin{split}
 \mathfrak G_n&=\{X:\ X\text{ is wall-free and shift-good}\},\\
 \mathfrak W_n&=\{X:\ k\le cD_N,\ X\text{ is not wall-free}\},\\
 \mathfrak C_n&=\{X:\ k\le cD_N,\ X\text{ is wall-free and a defect}\},\\
 \mathfrak L_n&=\{X:\ k>cD_N,\ X\notin\mathfrak G_n\}.
\end{split}
\end{equation}
The first class allows every size.  These four classes cover every
support exactly once.  This terminal size split is separate from the
fixed local-map threshold \(n_{\rm sm}\).  A support in
\(\mathfrak W_n\) is automatically shift-good when \(c\) is small:
if its padded rectangle also met the unit, its occupied extrema in
a wall-normal coordinate would span a fixed multiple of \(n\) minus
\(O(m')\).  Fixed-distance connectivity would then force
\(k\ge c_0n/m'\), contradicting \(k\le cD_N\).  This argument uses
no occupied unit block in an unlabelled defect.

Let \(\mathcal I_n=\{B:\{B\}\in\mathfrak G_n\}\) be the singleton
blocks in the first class.  On such a block, the centered and unit
activities are the plane activities at \(\phi\) and
\(\phi+\lambda_n\).  The fundamental singleton vanishes under the
phase projection, and the energy singleton contributes exactly
\[
 \frac{t_{j'}}2
  \E\{e_B^0(\zeta_{*,n}+\lambda_n)-e_B^0(\zeta_{*,n})\}
                    =\frac{t_{j'}}2e_B^0(\lambda_n).
\]
The cross term vanishes because \(\zeta_{*,n}\) is centered.  On all
supports in \(\mathfrak G_n\), the plane remainder has norm
\(O(N^{-2})\).  Its single-activity sum is bounded by the full
rooted estimate \(O(D_N^2/N^2)\), with the weight reserve for the
lift.  This is an upper bound for the remainder on
\(\mathfrak G_n\); the other three classes are treated below.

To identify the energy left by these singletons, define an allocation
on every terminal block, including the omitted ones, directly from
the connection:
\begin{equation}\label{obs:endpoint-allocation}
 E_{n,B}=\frac1{2a}
   \sum_{\substack{x\in B,\ e\in\{\pm e_1,\pm e_2\}\\
                   \{x,x+e\}\text{ internal to the square}}}
                      \eta_n(x,e)^2.
\end{equation}
Each internal unoriented edge is counted at its two endpoints, so
\[
       \sum_B E_{n,B}=g_n/a,\qquad
       E_{n,B}=e_B^0(\lambda_n)\quad(B\in\mathcal I_n)
\]
by \eqref{eq:singletons}.  Thus the missing central allocations are
defined without evaluating a lift at the unit.

For late \(N\), the complement of \(\mathcal I_n\) is a fixed-width
strip of terminal blocks at the walls together with the same central
singleton set in the two centered grids.  The wall-strip energy in
each square is \(O(m'/n)=O(D_N^{-1})\) by
\eqref{ent:unit-decay}.  On the identical central set, the difference
of the two endpoint energy densities is bounded by
\[
 C\left(\frac{1}{n_0(1+|x|)}+\frac1{n_0^2}\right).
\]
Write the difference of squares using the plane unit and
\eqref{ent:unit-decay} to obtain this bound.  Its sum over
\(|x|\le Cm'\), with the same internal neighbor stencils, is
\(O(m'/n_0)=O(D_N^{-1})\).  Consequently
\begin{equation}\label{obs:energy-allocation}
 \left[\sum_{B\in\mathcal I_n}e_B^0(\lambda_n)\right]_{n_0}^{n_1}
       =\frac{g_{n_1}-g_{n_0}}a+O(D_N^{-1}).
\end{equation}
Only this two-size difference is asserted; neither omitted central
energy is required to be small by itself.  Since \(t_{j'}=O(N^{-1})\),
the allocation error contributes \(O((ND_N)^{-1})\) to
\eqref{obs:terminal-energy}.

For \(X\in\mathfrak W_n\), its sites remain
a fixed fraction of \(n\) from the unit.  The first two difference
bounds in \eqref{ent:unit-decay}, and paths of \(O(k)\) blocks after
one anchor is removed, give
\[
                 \|\lambda_n\|_{m',D_X}\le C(1+k)/D_N.
\]
The neutral projection of the untwisted wall activity is even under
field negation.  Gaussian symmetry therefore makes the first
derivative at zero of
\(\E F_0(\zeta_{*,n}+r\lambda_n)\) vanish.  Taylor's formula, the
full analytic derivative norm, and the moment reserve in
Lemma~\ref{obs:terminal-family} bound its translated expectation
difference by
\[
             \frac{(1+k)^2 C^k}{N A_w^k D_N^2}.
\]
There are \(O_L(D_N)\) roots near a wall.  The rooted shape sum is
finite after the fixed weight choice, so the total contribution of
\(\mathfrak W_n\) is \(O((ND_N)^{-1})\).

For \(X\in\mathfrak C_n\), an individual activity can have norm
\(O(N^{-1})\) and need not be even.  A size-\(k\) defect has an
occupied block within \(O(k+1)\) block distance of the unit: its
padded rectangle reaches the unit and its block graph is connected
within a fixed distance.  Choose \(c\) small enough that every such
full padded set, including its stencils, lies in a common central
fraction of the smaller square.  These supports are then the same
in the two grids.
Their unit activities are the common plane prescription at
\(\phi+\delta_n\) by \eqref{obs:smooth-background}.  Modulo a common
constant,
\[
                 \|\delta_n\|_{m',D_X}\le C(1+k)/D_N.
\]

We compare the terminal fields on this full padded set.  Choose a
binary side \(v=2^b m'\asymp c_2n_0\), with fixed \(c_2>0\) small
enough that the polynomial range through \(v\), added to the set
and all its stencils, remains inside both squares.  On nonconstant
modes the exact split is
\[
 A_n^+\mathsf P_{m'}^n
 =\sum_{\substack{t=m',\,2m',\,\ldots\\t<v}}
       A_n^{-1}(\mathsf P_t^n-\mathsf P_{2t}^n)
       +A_n^+\mathsf P_v^n.
\]
The quotients in the sum use their polynomial continuations at
zero.  Independent Gaussian constants may be added to realize
these continuations because \(F_0\) is constant-invariant.
Finite range makes the restrictions of every polynomial piece to
the full padded set identical covariance matrices in the two
squares.  For this expectation couple those restrictions identically.
The unmatched tails \(\zeta_{{\rm long},n}\) satisfy, for every
fixed \(p<\infty\) and \(1\le i\le4\),
\[
 \sup_x\|\nabla^i\zeta_{{\rm long},n}(x)\|_{L^p}
                                      \le C_{p,i}n_0^{-i}.
\]
This is the positive-order multiplier estimate in the proof of
Lemma~\ref{ent:terminal-regulator}, now at \(v\asymp n_0\).
The same local discrete Sobolev estimate applied to
\(m'\nabla\zeta_{{\rm long},n}\) and
\((m')^2\nabla^2\zeta_{{\rm long},n}\), using two further differences,
gives their block maxima with bounds \(C_p/D_N\) and \(C_p/D_N^2\).
Paths in an enclosing rectangle of \(O((1+k)^2)\) blocks control
values relative to one anchor.  Hence the two fields with their
backgrounds, after subtracting their respective anchor values, obey
\begin{equation}\label{obs:coupling-direction}
 \left\|\,\|\xi_1-\xi_0\|_{m',D_X}\,\right\|_{L^p}
                      \le C_p(1+k)^C/D_N.
\end{equation}
Here \(\xi_i\) is the restriction of
\(\zeta_{*,n_i}+\delta_{n_i}\) in the coupling, with its anchor
removed.  The identical short fields cancel on the padded set.

The point norms at these two endpoints have the same local regulator
forms.  To see this precisely, let \(M_{D_X}\) be the positive matrix
representing the open energy \(E_{D_X}^0\) on its observation
coordinates, and let \(L_{D_X}\) be the joint Hilbert map from the
old Gaussian and noise coordinates to
\(B_s\alpha\zeta+e\) on those coordinates.  The internal edges and
whole observation cells are identical on the two padded sets, even
when those sets have holes or several components.  Thus
\(M_{D_X}\) is identical.  The restriction of
\(C_{<m'}=A^{-1}(I-\mathsf P_{m'})\) is also identical by its
polynomial finite range and the wall clearance; the noise law is
fixed.  Thus \(L_{D_X}L_{D_X}^*\) agrees.  The optimized formula
\cite[Equation~(9.16)]{B} depends on the Hilbert map only through
this covariance, as its convergent inverse series shows.  Hence
\(V_{m',h_{\rm b}}\) agrees, and the ordinary regulator has the
same edges and stencils as well.  This equality is for the local
high-mode regulator.  The terminal covariances are the different
covariances just coupled.

Remove the anchor pointwise from the neutral activity and apply the
mean value theorem on the segment from \(\xi_0\) to \(\xi_1\).
Use \eqref{obs:coupling-direction} in the exponent conjugate to
\(p_{\rm seg}\).  The identical local forms and the endpoint moment
\eqref{obs:terminal-moment}, with the deterministic backgrounds
\(\delta_n\), bound the regulator at each point of this segment.
For this support the expectation difference is at most
\[
                         \frac{(1+k)^C C^k}{ND_NA_w^k}.
\]
The same comparison applies to the centered activity, with zero
background.  The location bound for a size-\(k\) defect gives only
polynomially many possible roots.  The rooted shape and weight sum
absorbs these factors, so \(\mathfrak C_n\) contributes
\(O((ND_N)^{-1})\) in total.

Finally the supports in \(\mathfrak L_n\) have \(k>cD_N\).  The
absolute rooted bound with the remaining block weight gives
\[
               C D_N^2N^{-1}e^{-c'D_N}=O((ND_N)^{-1}).
\]
This includes the large supports reaching both a wall and the
unit.  Large supports in \(\mathfrak G_n\) were already included
in its plane remainder estimate.

The partition has now accounted for every single activity: the
plane energy and the vanishing fundamental on \(\mathcal I_n\),
the plane remainder on \(\mathfrak G_n\), and all terms in
\(\mathfrak W_n,\mathfrak C_n,\mathfrak L_n\).  Families with at least
two constituents contribute the separate \(O(D_N^4/N^2)\) error
from \eqref{obs:terminal-expansion}.  Combining these estimates
with \eqref{obs:energy-allocation} proves
\eqref{obs:terminal-energy}.  Substitution in
\eqref{obs:after-cancellation} proves
\eqref{obs:consecutive-formula}, absorbing the exponentially small
error.  The reference-energy calculation proved
\eqref{obs:green-constant}, completing Proposition~\ref{obs:consecutive}.
